% =====================================================================
% Preamble for eulerian_formal.tex
% A Stanley-style mathematical chapter with formal Rocq inline.
% =====================================================================

\documentclass[11pt,a4paper]{report}

% --- Encoding & fonts ---
\usepackage[T1]{fontenc}
\usepackage[utf8]{inputenc}
\usepackage{lmodern}
\usepackage{microtype}

% --- Math ---
\usepackage{amsmath, amssymb, amsthm, mathtools}
\usepackage{stmaryrd}  % \llbracket \rrbracket

% --- Page geometry ---
\usepackage[margin=1in]{geometry}

% --- Hyperlinks ---
\usepackage[colorlinks=true,
            linkcolor=blue!50!black,
            citecolor=blue!50!black,
            urlcolor=blue!50!black]{hyperref}
\usepackage{bookmark}

% --- Tables ---
\usepackage{booktabs}

% --- Code listings ---
\usepackage{xcolor}
\definecolor{rocqkw}{RGB}{120,30,150}
\definecolor{rocqcom}{RGB}{0,128,0}
\definecolor{rocqstr}{RGB}{200,30,30}
\definecolor{rocqbg}{RGB}{248,248,248}
\definecolor{rocqframe}{RGB}{220,220,230}
\definecolor{decodedbg}{RGB}{255,253,238}
\definecolor{sidebarbg}{RGB}{240,245,255}
\definecolor{sidebarframe}{RGB}{180,200,230}
\definecolor{computebg}{RGB}{245,255,245}

\usepackage{listings}

\lstdefinelanguage{Rocq}{
  morekeywords={Definition, Theorem, Lemma, Corollary, Proposition,
    Proof, Qed, Defined, Admitted, Fixpoint, Inductive, Variable,
    Variables, Hypothesis, Section, End, Module, Notation, Implicit,
    Types, Set, Unset, Arguments, From, Require, Import, Export,
    forall, exists, fun, match, with, end, if, then, else, let, in,
    return, as, structural, fix, cofix, struct, Type, Prop, Set,
    bool, nat, true, false, Some, None, option, list, prod, sum,
    Compute, Print, Check, Eval, Goal, Show, About, Reset, Open,
    Scope, Local},
  morekeywords=[2]{descent_set, des, asc, eulerian, beta, is_descent,
    descent, perm, set, ord_max, ord0, widen_ord, lift, val,
    insert_max_perm, insert_max_fun, extract_max_perm, ssreflect,
    fingroup},
  sensitive=true,
  morecomment=[s]{(*}{*)},
  morestring=[b]",
  basicstyle=\ttfamily\small,
  keywordstyle=\color{rocqkw}\bfseries,
  keywordstyle=[2]\color{rocqkw}\bfseries,
  commentstyle=\color{rocqcom}\itshape,
  stringstyle=\color{rocqstr},
  backgroundcolor=\color{rocqbg},
  frame=single,
  framerule=0.5pt,
  rulecolor=\color{rocqframe},
  framesep=4pt,
  xleftmargin=10pt,
  xrightmargin=10pt,
  breaklines=true,
  breakatwhitespace=true,
  showstringspaces=false,
  % UTF-8 chars that genuinely appear in our .v files (e.g. sigma in
  % qfact.v / qeul.v).  Without these, listings' default 8-bit handler
  % rejects the bytes with "Invalid UTF-8" errors.
  extendedchars=true,
  inputencoding=utf8,
  literate={->}{{$\to$}}{1}
           {=>}{{$\Rightarrow$}}{1}
           {>=}{{$\geq$}}{1}
           {<=}{{$\leq$}}{1}
           {<>}{{$\neq$}}{1}
           {forall}{{$\forall$}}{1}
           {exists}{{$\exists$}}{1}
           {σ}{{$\sigma$}}{1}
           {τ}{{$\tau$}}{1}
           {α}{{$\alpha$}}{1}
           {β}{{$\beta$}}{1}
           {γ}{{$\gamma$}}{1}
           {δ}{{$\delta$}}{1}
           {ε}{{$\varepsilon$}}{1}
           {π}{{$\pi$}}{1}
           {ω}{{$\omega$}}{1}
           {Σ}{{$\Sigma$}}{1}
           {ψ}{{$\psi$}}{1}
           {φ}{{$\phi$}}{1}
}

% --- (no lstnewenvironment wrapper; rocqbox uses lstlisting directly) ---

% --- Math notation aligned with Stanley ---
\newcommand{\Sn}[1]{\mathfrak{S}_{#1}}      % Stanley's S_n
\newcommand{\Snp}[1]{\mathfrak{S}_{#1}}     % alias
\newcommand{\des}{\mathrm{des}}
\newcommand{\asc}{\mathrm{asc}}
\newcommand{\maj}{\mathrm{maj}}
\newcommand{\inv}{\mathrm{inv}}
\newcommand{\descent}{\mathrm{D}}
\newcommand{\ord}[1]{\mathtt{'I_{#1}}}      % 'I_n
\newcommand{\Set}[1]{\mathtt{\{set\ #1\}}}
\newcommand{\Perm}[1]{\mathtt{\{perm\ #1\}}}
\newcommand{\eulerianA}{A}                  % Stanley's Eulerian A(n,k)
\newcommand{\Eul}{E}                        % Euler numbers
\newcommand{\beps}[2]{\beta_{#1}(#2)}       % Stanley's β_n(S)
\newcommand{\stirling}[2]{c(#1, #2)}        % Stanley's c(n,k)

% --- Theorem environments (Stanley-like) ---
\theoremstyle{plain}
\newtheorem{theorem}{Theorem}[chapter]
\newtheorem{proposition}[theorem]{Proposition}
\newtheorem{lemma}[theorem]{Lemma}
\newtheorem{corollary}[theorem]{Corollary}

\theoremstyle{definition}
\newtheorem{definition}[theorem]{Definition}
\newtheorem{example}[theorem]{Example}
\newtheorem{remark}[theorem]{Remark}

% --- Custom: formal/informal pairing ---
%
% \rocqsource{file}{line}  — clickable [Rocq, file:line] link to GitHub.
% Uses \GitHubBase and \GitHubBranch from below.

\newcommand{\GitHubBase}{https://github.com/LLM4Rocq/mathcomp-eulerian/blob}
\newcommand{\GitHubBranch}{main}

% \rocqsource takes a filename which may contain underscores.
% \detokenize keeps the underscores literal in BOTH the URL argument
% of \href AND the visible text (otherwise `_` is the subscript catcode
% in body text and triggers a math-mode error).
\newcommand{\rocqsource}[2]{%
  \href{\GitHubBase/\GitHubBranch/\detokenize{#1}\#L#2}%
       {\textsf{[\,Rocq, \texttt{\detokenize{#1}:#2}\,]}}%
}

\newcommand{\rocqsourcerange}[3]{%
  \href{\GitHubBase/\GitHubBranch/\detokenize{#1}\#L#2-L#3}%
       {\textsf{[\,Rocq, \texttt{\detokenize{#1}:#2-#3}\,]}}%
}

% Usage pattern:
%
%   \rocqsource{file.v}{LINE}
%   \begin{lstlisting}[language=Rocq]
%   ...formal Rocq text verbatim...
%   \end{lstlisting}
%
% (Wrapping listings in \newenvironment is fragile across multi-pass
%  builds with hyperref, so we use the listing environment directly.)

% --- Pedagogical sidebars ---

% \decoded{Plain-text mechanical translation of the formal block.}
% Appears immediately after a Rocq block; substitutes mathcomp casts for
% their meaning so the reader can see the formal text says what the
% informal text claims.
\newcommand{\decoded}[1]{%
  \par\nopagebreak\smallskip%
  \noindent\colorbox{decodedbg}{%
    \begin{minipage}{\dimexpr\textwidth-2\fboxsep\relax}%
      \footnotesize\textsf{Decoded.}\enspace #1%
    \end{minipage}%
  }\par\smallskip%
}

% \begin{sidebar}{Title} ... \end{sidebar}
% A "Why this representation?" or "Design choice" box, used to explain
% formalization decisions inline.  Implemented via a quote-style indent
% with a coloured rule on the left margin.
\newenvironment{sidebar}[1]{%
  \par\smallskip%
  \noindent\hspace{0.02\textwidth}%
  \begin{minipage}{0.96\textwidth}%
    \small%
    {\color{sidebarframe}\rule{2pt}{0.6\baselineskip}}\enspace%
    \textbf{Design choice: #1}\par\smallskip%
}{%
  \end{minipage}\par\smallskip%
}

% \begin{computestrip} ... \end{computestrip}
% A short Rocq Compute output strip, demonstrating a definition's value
% on a concrete example.
\lstnewenvironment{computestrip}{%
  \lstset{language=Rocq, backgroundcolor=\color{computebg}}%
}{}

% --- Title ---
\title{The mathcomp-eulerian formalization\\
  \large A Stanley-style reading of permutation statistics in Rocq/MathComp}
\author{}
\date{}


\begin{document}

\maketitle

\tableofcontents

% =====================================================================
\chapter*{About this document}
% =====================================================================

\addcontentsline{toc}{chapter}{About this document}

This document is a Stanley-style reading of the
\href{https://github.com/LLM4Rocq/mathcomp-eulerian}{\texttt{mathcomp-eulerian}}
formalization. Each definition and statement is given twice: once in
ordinary mathematical English (in the spirit of Stanley
\emph{Enumerative Combinatorics I}, Chapter~1) and once as the
verbatim Rocq text, with a clickable link to the line in the
repository on GitHub. The intention is that a mathematician can read
top to bottom like a textbook and at any point click through to
verify the formal statement.

\paragraph{Status of this prototype.} This first edition contains
only Chapter~\ref{ch:descents} (Stanley \S1.4: descents, Eulerian
numbers, and the set-refined count $\beps{n}{S}$). Chapters covering
\S1.3 (cycles, inversions, Foata), the $q$-analogues, \S1.6.2
(longest alternating subsequence), and \S1.6.3 (toggle action and
the headline corollary) are planned but not yet written.

\paragraph{Conventions.} The single recurring shift is 0-indexing.
Stanley writes $w \in \Sn{n}$ for permutations of $\{1, \dots, n\}$
and $\descent(w) \subseteq \{1, \dots, n-1\}$ for the descent set;
we use $\Perm{\ord{n}}$ for permutations of
$\{0, \dots, n-1\}$ and $\Set{\ord{n-1}}$ for the descent set.
Concretely:
\[
  \text{our } \Perm{\ord{n+1}} \;=\; \text{Stanley's } \Sn{n+1}.
\]
A complete translation table is in
\href{https://github.com/LLM4Rocq/mathcomp-eulerian/blob/main/docs/READING_GUIDE.md}{\texttt{docs/READING\_GUIDE.md}};
a worked example through Stanley's $w = [3,1,4,2]$ is in the same file.

\paragraph{Reading the formal text.} Each Rocq block looks like
this:

\rocqsource{descent.v}{25}
\begin{lstlisting}[language=Rocq]
Definition descent_set s : {set 'I_n} := [set i | is_descent s i].
\end{lstlisting}

\noindent The bracketed link at the top is clickable and opens the
exact line on GitHub. The body is the Rocq source verbatim;
identifiers and keywords are highlighted but no mathematical
substitution has been performed. If the formal block looks different
from what you expected from the informal statement above it, that is
exactly the kind of discrepancy worth flagging.

\paragraph{What is \emph{not} included here.}
Full proofs are not included; only the formal \emph{statements} are
shown alongside the informal mathematics, with a one-paragraph proof
sketch for each non-trivial result. The full proof is always one
click away in the linked source. The cd-index development
(\texttt{mmtree.v}, \texttt{psi\_*.v}) is project-internal scaffolding
not in Stanley and is not covered.

\paragraph{Verifying axiom-freeness.} Every theorem in this document
is kernel-validated by Rocq with no axioms or admits. To check this
yourself for, e.g., the headline Corollary~\ref{cor:beta_alt_max}:

\begin{verbatim}
echo 'From mathcomp_eulerian Require Import beta_swap.
      Print Assumptions beta_alt_max.' | rocq top -R . mathcomp_eulerian
\end{verbatim}

\noindent prints \texttt{Closed under the global context}, meaning no
axiom is used anywhere in its transitive dependencies.

\clearpage

% =====================================================================
\chapter{Descents and Eulerian numbers}
\label{ch:descents}
% =====================================================================

This chapter formalizes Stanley \S1.4 (Stanley
\emph{Enumerative Combinatorics I}, 2nd ed., pp.~30--38). The central
objects are the \emph{descent set} of a permutation, the
\emph{Eulerian numbers} $\eulerianA(n, k)$ that count permutations by
descent count, and the set-refined count $\beps{n}{S}$ that counts
permutations by their full descent set.

% =====================================================================
\section{The descent set}
\label{sec:descent_set}
% =====================================================================

\begin{definition}[Descent]
  \label{def:is_descent}
  Let $w = (w_1, w_2, \dots, w_n) \in \Sn{n}$. A position
  $i \in \{1, \dots, n-1\}$ is a \emph{descent} of $w$ if
  $w_i > w_{i+1}$.
\end{definition}

\rocqsource{descent.v}{22}
\begin{lstlisting}[language=Rocq]
Definition is_descent s i : bool :=
  s (widen_ord (leqnSn n) i) > s (lift ord0 i).
\end{lstlisting}

\noindent
Here $s$ has type $\Perm{\ord{n+1}}$ and $i$ has type $\ord{n}$.
The cast \texttt{widen\_ord (leqnSn n) i} sends $i : \ord{n}$ to
the same numerical value in $\ord{n+1}$ (the position $i$ itself);
the cast \texttt{lift ord0 i} sends $i$ to $i+1 : \ord{n+1}$. Hence
the formal definition reads $s(i) > s(i+1)$, matching
Definition~\ref{def:is_descent} modulo 0-indexing.

\begin{definition}[Descent set]
  \label{def:descent_set}
  The \emph{descent set} of $w \in \Sn{n}$ is
  \[
    \descent(w) \;=\; \{\, i \in \{1, \dots, n-1\} : w_i > w_{i+1} \,\}.
  \]
\end{definition}

\rocqsource{descent.v}{25}
\begin{lstlisting}[language=Rocq]
Definition descent_set s : {set 'I_n} := [set i | is_descent s i].
\end{lstlisting}

\noindent
$\descent(w)$ is a subset of $\{1, \dots, n-1\}$ in Stanley; in our
types it is a subset of $\{0, \dots, n-1\} = \ord{n}$. The
correspondence is $\text{Stanley's } i \leftrightarrow \text{our } (i-1)$.

\begin{example}
  \label{ex:s3142}
  Take $w = [3, 1, 4, 2] \in \Sn{4}$. Then
  $\descent(w) = \{1, 3\}$ since $w_1 = 3 > 1 = w_2$ and
  $w_3 = 4 > 2 = w_4$. In our types, the corresponding $s :
  \Perm{\ord{4}}$ has $s\,0 = 2$, $s\,1 = 0$, $s\,2 = 3$, $s\,3 = 1$
  (values shifted to $\{0,1,2,3\}$), and $\mathtt{descent\_set\ s} =
  \{0, 2\} : \Set{\ord{3}}$.
\end{example}

\begin{definition}[Descent number, ascent number]
  \label{def:des_asc}
  $\des(w) = |\descent(w)|$ and $\asc(w) = (n-1) - \des(w)$.
\end{definition}

\rocqsource{descent.v}{27}
\begin{lstlisting}[language=Rocq]
Definition des s : nat := #|descent_set s|.
\end{lstlisting}

\rocqsource{descent.v}{29}
\begin{lstlisting}[language=Rocq]
Definition asc s : nat := n - des s.
\end{lstlisting}

\noindent
In the formal definition, $n$ is the number of descent positions
(Stanley's $n - 1$). So our $\mathtt{asc\ s} = n - \mathtt{des\ s}$
matches Stanley's $(n-1) - \des(w)$.

\begin{lemma}[Descent count is bounded]
  \label{lem:des_le}
  $\des(w) \leq n - 1$ for every $w \in \Sn{n}$.
\end{lemma}

\rocqsource{descent.v}{38}
\begin{lstlisting}[language=Rocq]
Lemma des_le s : des s <= n.
\end{lstlisting}

\noindent
\emph{Proof sketch.} $\descent(w) \subseteq \{0, \dots, n-1\}$ which
has $n$ elements; the cardinality is at most $n$. \qed

\begin{lemma}[Identity has no descents]
  \label{lem:des_id}
  $\des(\mathrm{id}) = 0$.
\end{lemma}

\rocqsource{descent.v}{44}
\begin{lstlisting}[language=Rocq]
Lemma des_id : des (1 : {perm 'I_n.+1}) = 0.
\end{lstlisting}

\noindent
\emph{Proof sketch.} For the identity, $s\,i = i$ at every position,
so $s(i) = i < i+1 = s(i+1)$ --- no descent slot. \qed

% =====================================================================
\section{Reversal symmetry}
\label{sec:reversal}
% =====================================================================

A useful symmetry: reversing a permutation flips ascents and
descents at corresponding positions.

\begin{definition}[Reverse permutation]
  \label{def:rev_perm}
  Let $w \in \Sn{n}$ and define $w^{\mathrm{rev}}$ by
  $w^{\mathrm{rev}}_i = w_{n+1-i}$.
\end{definition}

\rocqsource{descent.v}{80}
\begin{lstlisting}[language=Rocq]
Definition rev_perm_ord : {perm 'I_n.+1} := perm (@rev_ord_inj n.+1).
\end{lstlisting}

\rocqsource{descent.v}{85}
\begin{lstlisting}[language=Rocq]
Definition rev_perm s : {perm 'I_n.+1} := rev_perm_ord * s.
\end{lstlisting}

\noindent
\texttt{rev\_perm\_ord} is the position-reversal involution
$i \mapsto n - i$; \texttt{rev\_perm s} pre-composes $s$ with it,
giving Stanley's $w^{\mathrm{rev}}$.

\begin{lemma}[Descents of the reverse]
  \label{lem:is_descent_rev}
  Position $i$ is a descent of $w^{\mathrm{rev}}$ iff position $n - i$
  is \emph{not} a descent of $w$:
  \[
    i \in \descent(w^{\mathrm{rev}})
    \iff
    (n - i) \notin \descent(w).
  \]
\end{lemma}

\rocqsource{descent.v}{90}
\begin{lstlisting}[language=Rocq]
Lemma is_descent_rev s (i : 'I_n) :
  is_descent (rev_perm s) i = ~~ is_descent s (rev_ord i).
\end{lstlisting}

\noindent
\emph{Proof sketch.} Direct calculation: $\mathtt{(rev\_perm\,s)\,j} =
s(n-j)$, and the comparison $s(n-i) > s(n-i-1)$ is the negation of
$s(n-i-1) > s(n-i)$. \qed

% =====================================================================
\section{Eulerian numbers}
\label{sec:eulerian}
% =====================================================================

\begin{definition}[Eulerian number]
  \label{def:eulerian}
  The \emph{Eulerian number} $\eulerianA(n, k)$ is the number of
  permutations of $\{1, \dots, n\}$ with exactly $k$ descents:
  \[
    \eulerianA(n, k) \;=\; \#\{\, w \in \Sn{n} : \des(w) = k \,\}.
  \]
\end{definition}

\rocqsource{eulerian.v}{14}
\begin{lstlisting}[language=Rocq]
Definition eulerian (n k : nat) : nat :=
  #|[set s : {perm 'I_n.+1} | des s == k]|.
\end{lstlisting}

\noindent
\textbf{Off-by-one alert.} The Rocq parameter $n$ is "permutation
size minus one." Concretely, \texttt{eulerian n k} counts perms of
$\Perm{\ord{n+1}}$, which are Stanley's $\Sn{n+1}$. Hence
\[
  \texttt{eulerian } n\ k \;=\; \eulerianA(n+1, k).
\]
For instance, \texttt{eulerian 3 1} corresponds to Stanley's
$\eulerianA(4, 1) = 11$ (Stanley table p.~36).

\begin{lemma}[Row sum]
  \label{lem:eulerian_row_sum}
  $\sum_{k=0}^{n-1} \eulerianA(n, k) = n!$
\end{lemma}

\rocqsource{eulerian.v}{31}
\begin{lstlisting}[language=Rocq]
Lemma eulerian_row_sum_fact n :
  \sum_(k < n.+1) eulerian n k = n.+1`!.
\end{lstlisting}

\noindent
\emph{Proof sketch.} The summands partition $\Sn{n+1}$ by descent
count; the total is $|\Sn{n+1}| = (n+1)!$. \qed

\begin{lemma}[Out-of-range vanishing]
  \label{lem:eulerian_out_of_range}
  $\eulerianA(n, k) = 0$ for $k \geq n$.
\end{lemma}

\rocqsource{eulerian.v}{34}
\begin{lstlisting}[language=Rocq]
Lemma eulerian_out_of_range n k :
  n < k -> eulerian n k = 0.
\end{lstlisting}

\noindent
\emph{Proof sketch.} A permutation has at most $n$ descent positions
(Lemma~\ref{lem:des_le}); none can have $k > n$ descents. \qed

\begin{lemma}[The unique perm with no descents is the identity]
  \label{lem:des0_id}
  If $\des(w) = 0$ then $w = \mathrm{id}$.
\end{lemma}

\rocqsource{eulerian.v}{41}
\begin{lstlisting}[language=Rocq]
Lemma des0_id n (s : {perm 'I_n.+1}) : des s = 0 -> s = 1%g.
\end{lstlisting}

\noindent
\emph{Proof sketch.} $\des(s) = 0$ means $s(i) < s(i+1)$ for every
adjacent pair $i$, so $s$ is strictly increasing. The only strictly
increasing function $\ord{n+1} \to \ord{n+1}$ is the identity. \qed

% =====================================================================
\section{The insert-max bijection}
\label{sec:insert_max}
% =====================================================================

A central tool for the Eulerian recurrence is the bijection that
inserts the maximum value $n+1$ at a chosen position into a
permutation of $\{1, \dots, n\}$.

\begin{definition}[Insert-max]
  \label{def:insert_max}
  For $t \in \Sn{n}$ and $p \in \{1, \dots, n+1\}$, let
  $\mathtt{insert\_max}(t, p) \in \Sn{n+1}$ be the permutation
  obtained by inserting the value $n+1$ at position $p$ in the
  one-line form of $t$.
\end{definition}

\rocqsourcerange{eulerian.v}{134}{157}
\begin{lstlisting}[language=Rocq]
Definition insert_max_fun (i : 'I_n.+2) : 'I_n.+2 :=
  match unlift p i with
  | Some j => widen_ord (leqnSn _) (t j)
  | None => ord_max
  end.

Definition insert_max_perm : {perm 'I_n.+2} := perm insert_max_fun_inj.
\end{lstlisting}

\noindent
At position $p$, the output is $\mathtt{ord\_max} = n+1$. At any
other position $\mathtt{lift\ p\ j}$ (i.e., positions $\neq p$,
parameterized by $j : \ord{n+1}$), the output is
$t\,j$ cast back to $\ord{n+2}$.

\begin{definition}[Extract-max]
  \label{def:extract_max}
  The inverse: given $\sigma \in \Sn{n+1}$ with $\sigma(p) = n+1$
  for some position $p$, remove that value and standardise to obtain
  $t \in \Sn{n}$.
\end{definition}

\rocqsourcerange{eulerian.v}{185}{203}
\begin{lstlisting}[language=Rocq]
Definition extract_max_fun (j : 'I_n.+1) : 'I_n.+1 :=
  odflt j (unlift ord_max (s (lift p j))).

Definition extract_max_perm : {perm 'I_n.+1} := perm extract_max_fun_inj.
\end{lstlisting}

\begin{theorem}[Insert-max is a bijection]
  \label{thm:insert_max_bij}
  The map $(t, p) \mapsto \mathtt{insert\_max}(t, p)$ is a bijection
  $\Sn{n} \times \{1, \dots, n+1\} \to \Sn{n+1}$.
\end{theorem}

\rocqsource{eulerian.v}{437}
\begin{lstlisting}[language=Rocq]
Lemma insert_max_perm_bij n :
  forall sigma : {perm 'I_n.+2},
    exists! tp : {perm 'I_n.+1} * 'I_n.+2,
      sigma = insert_max_perm tp.1 tp.2.
\end{lstlisting}

\noindent
\emph{Proof sketch.} Given $\sigma$, the unique $p$ is
$\sigma^{-1}(n+1)$ and $t = \mathtt{extract\_max}(\sigma, p)$. The
two roundtrips are direct calculations. \qed

% =====================================================================
\section{Eulerian recurrence}
\label{sec:eulerian_rec}
% =====================================================================

\begin{theorem}[Eulerian recurrence]
  \label{thm:eulerian_rec}
  For $n \geq 1$ and $0 \leq k \leq n$,
  \[
    \eulerianA(n+1, k)
    \;=\;
    (k+1)\,\eulerianA(n, k) \;+\; (n - k + 1)\,\eulerianA(n, k-1).
  \]
\end{theorem}

\rocqsource{eulerian.v}{458}
\begin{lstlisting}[language=Rocq]
Theorem eulerian_rec n k :
  eulerian n.+1 k =
    (k.+1) * eulerian n k + (n.+1 - k) * eulerian n k.-1.
\end{lstlisting}

\noindent
\emph{Proof sketch.} Using the insert-max bijection, classify each
$\sigma \in \Sn{n+1}$ by the position $p$ of $n+1$. There are
$k+1$ ascent positions and $n - k$ descent positions where inserting
$n+1$ preserves the descent count; inserting at the other $n - k + 1$
slots increases it by one. Summing recovers the recurrence. \qed

% =====================================================================
\section{The set-refined count $\beta$}
\label{sec:beta}
% =====================================================================

The descent count is a coarse invariant. The full descent
\emph{set} carries more information; the count of permutations with
a fixed descent set is a finer object.

\begin{definition}[$\beta$ count]
  \label{def:beta}
  For a subset $S \subseteq \{1, \dots, n-1\}$, the count
  $\beps{n}{S}$ is the number of permutations of $\{1, \dots, n\}$
  whose descent set is exactly $S$:
  \[
    \beps{n}{S} \;=\; \#\{\, w \in \Sn{n} : \descent(w) = S \,\}.
  \]
\end{definition}

\rocqsource{beta.v}{19}
\begin{lstlisting}[language=Rocq]
Definition beta (n : nat) (D : {set 'I_n}) : nat :=
  #|[set sigma : {perm 'I_n.+1} | descent_set sigma == D]|.
\end{lstlisting}

\noindent
\textbf{Off-by-one alert.} As with \texttt{eulerian}, the Rocq
parameter $n$ is "perm size minus one." So
\texttt{beta n D} for $D : \Set{\ord{n}}$ equals Stanley's
$\beps{n+1}{D}$, with $D$ interpreted as a subset of $\{0, \dots, n-1\}$
that corresponds to Stanley's 1-indexed $S \subseteq \{1, \dots, n\}$.

\begin{lemma}[$\beta$ partitions Eulerian]
  \label{lem:beta_eulerian}
  Summing $\beps{n}{S}$ over all $S$ of cardinality $k$ recovers the
  Eulerian number:
  \[
    \sum_{|S| = k} \beps{n}{S} \;=\; \eulerianA(n, k).
  \]
\end{lemma}

\rocqsource{beta.v}{124}
\begin{lstlisting}[language=Rocq]
Lemma beta_eulerian n k :
  \sum_(D : {set 'I_n} | #|D| == k) beta D = eulerian n k.
\end{lstlisting}

\noindent
\emph{Proof sketch.} The sets $\{\sigma : \descent(\sigma) = D\}$ for
$D$ ranging over subsets of cardinality $k$ partition the set of
permutations with $k$ descents. The cardinalities sum. \qed

\begin{lemma}[$\beta$ is reversal-invariant]
  \label{lem:beta_rev}
  For $S \subseteq \{1, \dots, n-1\}$, let $S^{\mathrm{rev}} =
  \{n - i : i \in S\}$. Then
  $\beps{n}{S^{\mathrm{rev}}} = \beps{n}{S^{\mathrm{c}}}$,
  where $S^{\mathrm{c}}$ is the complement in $\{1, \dots, n-1\}$.
\end{lemma}

\rocqsource{beta.v}{140}
\begin{lstlisting}[language=Rocq]
Lemma beta_rev n (D : {set 'I_n}) :
  beta [set rev_ord i | i in D] = beta (~: D).
\end{lstlisting}

\noindent
\emph{Proof sketch.} Reversal $w \mapsto w^{\mathrm{rev}}$ is a
bijection on $\Sn{n}$; by Lemma~\ref{lem:is_descent_rev} it maps
descent set $D$ to $\{n - i : i \notin D\}$, which is the
complement of $\{n - i : i \in D\}$. Counting on both sides
matches. \qed

% =====================================================================
\section{What this chapter doesn't cover}
% =====================================================================

This chapter establishes the basic objects ($\descent$, $\des$,
$\eulerianA$, $\beta$) and the foundational identities. The headline
of the formalization,
\[
  \beps{n}{S} \,\leq\, E_n
  \quad \text{with equality iff } S \text{ is alternating}
  \quad (\text{Stanley Cor 1.6.5})
\]
requires the toggle-action machinery of \S1.6.3 and is the subject
of a future chapter.

\appendix

% =====================================================================
\chapter{Source map}
% =====================================================================

A complete file index is in \texttt{docs/READING\_GUIDE.md} in the
\href{https://github.com/LLM4Rocq/mathcomp-eulerian}{\texttt{mathcomp-eulerian}}
repository. Files covered (or planned) by this document, in Stanley
reading order:

\begin{description}
  \item[\texttt{descent.v}, \texttt{eulerian.v}, \texttt{beta.v}]
    Stanley \S1.4 (this chapter).
  \item[\texttt{inversions.v}] Stanley \S1.3.3 (planned).
  \item[\texttt{cycles.v}] Stanley \S1.3.1--2 (planned).
  \item[\texttt{foata.v}] Stanley \S1.3.4 (planned).
  \item[\texttt{qfact.v}, \texttt{qeul.v}] Stanley \S1.4 $q$-analogues
    (planned).
  \item[\texttt{altsub.v}] Stanley \S1.6.2 (planned).
  \item[\texttt{beta\_omega.v}, \texttt{beta\_swap.v}] Stanley \S1.6.3
    (planned).
  \item[\texttt{experimental/reflection.v}] Stanley \S1.6.4
    (planned, contains 1 named admit).
\end{description}

\end{document}
